\documentclass[fontsize=15pt]{jlreq}

%プリアンブル
\usepackage{amsmath}
\usepackage{mathtools}
\pagestyle{empty}
\usepackage[top=15mm, bottom=20mm, left=20mm, right=20mm]{geometry}

\newcommand{\m}[1]{\raisebox{0.1em}{\textcircled{\scriptsize #1}}}



%本文
\begin{document}
\begin{center}
{\LARGE {振り返り}} \\
-物理\m{1}- %単元ごとに変えるか消す
\end{center}
\vspace{1em}
\begin{flushright}
名前：\underline{\hspace{5cm}}
\end{flushright}

% 問題部分
\begin{enumerate}
    \item 次の物理量の次元(単位)をSI単位系で書け．\\
    \begin{itemize}
        \item 変位・距離：
        \item 時間：
        \item 速度・速さ：
    \end{itemize}
    \vspace{2em} 

    \item 物体が5秒間で10[m]進んだ．この時の平均の速さを求めよ．\\
    (\hspace{4cm})
    \vspace{2em} 

    \item 平均の速さと瞬間の速さの違いについて説明せよ．\\
    (\hspace{15cm})
    
    \vspace{2em} 

    \item 物体が一直線上を一定の速さで進む運動をなんというか．\\
    (\hspace{4cm})
    \vspace{2em} 

    \item 1/50秒ごとに打点する記録タイマーがある．その記録タイマーが5回打点したとき，何秒経過したと言えるか．\\
    (\hspace{4cm})
    \vspace{2em}  

\end{enumerate}

エクストラ：ニュートンの三法則をそれぞれかけ
\begin{itemize}
    \item \m{1}
    \item \m{2}
    \item \m{3}
\end{itemize}

\end{document}
